Dem Profil des CAS entsprechend haben wir generative KI-Werkzeuge gezielt und transparent als Unterstützung eingesetzt. KI diente uns durchgängig als Ergänzung und Sparringpartner, nicht als Ersatz der eigenen Denk- und Gestaltungsleistung. Sämtliche KI-gestützten Zwischenergebnisse wurden von uns inhaltlich geprüft, kritisch eingeordnet und verantwortet. Die zentralen Erkenntnisse beruhen auf den selbst geführten Interviews und dem realen Prototypentest; KI-generierte Beiträge wurden bewusst nachrangig gewichtet.

\textbf{Eingesetzte KI-Werkzeuge und Verwendungszweck:}
\begin{itemize}
  \item \textbf{ChatGPT (OpenAI)} -- als Sparringpartner bei der Entwicklung des Interviewleitfadens, zur Strukturierung und sprachlichen Überarbeitung von Texten sowie für Reframing-Provokationen bei den How-might-we-Fragen.
  \item \textbf{KI-basierter Virtual Trainer} -- zum Üben und Reflektieren der Gesprächsführung vor den Leitfadeninterviews.
  \item \textbf{Synthetic Users} -- als KI-generierte Persona-Stellvertreter zur Verbreiterung der qualitativen Datenbasis; aufgrund ihrer Tendenz zu sozial erwünschten Antworten (\emph{People Pleasing}) bewusst nur ergänzend genutzt.
  \item \textbf{Generative Bildwerkzeuge} -- zur Erstellung des Titelbilds sowie der Plakat- und Kampagnenvisualisierung des Treffpunkts ``Donnerstags mit Kurt''.
\end{itemize}

\textbf{Kennzeichnung KI-generierter Bilder:} Das Titelbild dieses Berichts sowie die Plakatvisualisierung in der Develop-Phase wurden mit generativen Bildwerkzeugen erstellt. Alle übrigen Abbildungen sind eigene Darstellungen der Arbeitsgruppe und entsprechend als ``eigene Darstellung'' ausgewiesen.

Die Verantwortung für die inhaltliche Richtigkeit, die Auswahl der Quellen und die Schlussfolgerungen dieses Berichts liegt vollständig bei den Autor:innen.
